% unidades/17-planes-citas.tex
\chapter{📅 計画と予定 — Planes y Citas}
\unidadheader{17}{計画と予定}{Planes y Citas}{Hablar de intenciones, planes e invitaciones}

\begin{tcolorbox}[vocabbox, title={📌 Vocabulario Nuevo}]
\begin{tabularx}{\linewidth}{llX}
\toprule
\textbf{Japonés} & \textbf{Español} & \textbf{Tipo} \\
\midrule
\vocabitem{予定}{よてい}{plan / agenda}{sust.}
\vocabitem{計画}{けいかく}{plan / proyecto}{sust.}
\vocabitem{旅行}{りょこう}{viaje}{sust.}
\vocabitem{来週}{らいしゅう}{la semana que viene}{adv.}
\vocabitem{来月}{らいげつ}{el mes que viene}{adv.}
\vocabitem{来年}{らいねん}{el año que viene}{adv.}
\vocabitem{〜週間}{しゅうかん}{durante ~ semanas}{contador}
\vocabitem{〜ヶ月}{かげつ}{durante ~ meses}{contador}
\vocabitem{誘う}{さそう}{invitar / proponer}{v. gr.1}
\vocabitem{決める}{きめる}{decidir}{v. gr.2}
\bottomrule
\end{tabularx}
\end{tcolorbox}

\subsection*{🗣️ Frases Clave}

\frase{いっしょに\ruby{映画}{えいが}を\ruby{見}{み}ませんか？}
  {¿No quieres ver una película juntos?}
\frase{いいですね！\ruby{行}{い}きましょう！}{¡Qué buena idea! ¡Vamos!}
\frase{すみません、\ruby{予定}{よてい}があります。}
  {Lo siento, tengo planes.}
\frase{\ruby{来週}{らいしゅう}\ruby{京都}{きょうと}に\ruby{行}{い}くつもりです。}
  {Tengo la intención de ir a Kioto la semana que viene.}

\subsection*{⚙️ Gramática}

\patron{Intención: つもりです}
  {V-dict. + つもりです (tengo la intención de ~)}
  {\ej{\ruby{大学}{だいがく}に\ruby{入}{はい}るつもりです。}
    {Tengo la intención de entrar a la universidad.}
  \ej{たばこをやめるつもりです。}
    {Tengo la intención de dejar de fumar.}}

\patron{Plan / Agenda: 予定です}
  {V-dict. + 予定です (está planeado ~)}
  {\ej{\ruby{来月}{らいげつ}\ruby{日本}{にほん}に\ruby{行}{い}く\ruby{予定}{よてい}です。}
    {Está planeado que vaya a Japón el mes que viene.}}

\patron{Invitación: ませんか / ましょう}
  {V-ません + か (¿no quieres ~?) / V-ましょう (¡hagamos ~!)}
  {\ej{いっしょに\ruby{食}{た}べませんか？}{¿No comemos juntos?}
  \ej{そろそろ\ruby{行}{い}きましょう。}{Ya vámonos.}}

\begin{tcolorbox}[ejerciciobox, title={📝 Ejercicios Rápidos}]
\begin{enumerate}
  \item Invita a alguien a tomar café usando ませんか.
  \item Di tu plan para las vacaciones usando 予定です.
  \item ¿Cuál es la diferencia entre つもりです y 予定です?
\end{enumerate}
\vspace{0.4em}
\textbf{Respuestas:}
1) コーヒーを\ruby{飲}{の}みませんか？\quad
2) (personal)\quad
3) つもり = voluntad/intención personal; 予定 = plan ya establecido/agenda.
\end{tcolorbox}

\begin{tcolorbox}[kanjibox, title={🀄 Kanjis de la Unidad}]
\begin{tabularx}{\linewidth}{cllXX}
\toprule
\textbf{Kanji} & \textbf{On} & \textbf{Kun} & \textbf{Significado} & \textbf{Vocab} \\
\midrule
\kanjirow{予}{よ}{—}{previo}{\ruby{予定}{よてい} / \ruby{予約}{よやく}}
\kanjirow{定}{てい}{さだ}{fijar}{\ruby{予定}{よてい} / \ruby{決定}{けってい}}
\kanjirow{計}{けい}{はか}{medir / plan}{\ruby{計画}{けいかく} / \ruby{時計}{とけい}}
\kanjirow{画}{かく・が}{—}{trazo / cuadro}{\ruby{計画}{けいかく} / \ruby{映画}{えいが}}
\kanjirow{約}{やく}{—}{promesa / aprox.}{\ruby{予約}{よやく} / \ruby{約}{やく}}
\bottomrule
\end{tabularx}
\end{tcolorbox}
